% Decorative vector art (TikZ only, no raster), placed in reference-pixel coordinates.
\newcommand{\Pine}[4]{% x, base y, height, colour  -- three-tier pine
  \fill[#4] (#1,#2-#3) -- (#1-#3*0.16,#2-#3*0.55) -- (#1+#3*0.16,#2-#3*0.55) -- cycle;
  \fill[#4] (#1,#2-#3*0.78) -- (#1-#3*0.22,#2-#3*0.32) -- (#1+#3*0.22,#2-#3*0.32) -- cycle;
  \fill[#4] (#1,#2-#3*0.52) -- (#1-#3*0.27,#2-#3*0.1) -- (#1+#3*0.27,#2-#3*0.1) -- cycle;
  \fill[#4] (#1-#3*0.035,#2) rectangle (#1+#3*0.035,#2-#3*0.12);}
\newcommand{\Pylon}[4]{% x, base y, height, colour -- lattice transmission tower (line art)
  \draw[#4,line width=0.75pt,line join=round]
    (#1-#3*0.2,#2) -- (#1-#3*0.05,#2-#3*0.86) -- (#1,#2-#3) -- (#1+#3*0.05,#2-#3*0.86) -- (#1+#3*0.2,#2)
    (#1-#3*0.32,#2-#3*0.74) -- (#1+#3*0.32,#2-#3*0.74) (#1-#3*0.24,#2-#3*0.56) -- (#1+#3*0.24,#2-#3*0.56)
    (#1-#3*0.32,#2-#3*0.74) -- (#1-#3*0.32,#2-#3*0.68) (#1+#3*0.32,#2-#3*0.74) -- (#1+#3*0.32,#2-#3*0.68)
    (#1-#3*0.16,#2-#3*0.18) -- (#1+#3*0.1,#2-#3*0.48) (#1+#3*0.16,#2-#3*0.18) -- (#1-#3*0.1,#2-#3*0.48);}

% ---------------------------------------------------------------- header: faint hills (left) and city skyline (right)
\newcommand{\HillsLeft}{%
  \fill[PSky,opacity=0.7] (0,111) -- (0,86) .. controls (25,80) and (48,74) .. (74,77) .. controls (102,81) and (122,92) .. (150,90)
     .. controls (182,88) and (205,80) .. (232,86) .. controls (262,93) and (300,100) .. (345,111) -- cycle;
  \fill[PSky!60!PMist,opacity=0.45] (0,111) -- (0,99) .. controls (40,93) and (85,91) .. (130,98) .. controls (170,104) and (220,107) .. (265,111) -- cycle;}
% Header city skyline (right): low, light, behind the tagline, as in the reference (vector only).
\newcommand{\SkylineRight}{%
  % distant mountain range
  \fill[PSky,opacity=0.75] (930,111) .. controls (948,98) and (962,88) .. (980,83) .. controls (995,79) and (1010,72) .. (1032,70)
     .. controls (1055,68) and (1070,78) .. (1092,76) .. controls (1115,74) and (1135,70) .. (1161,72) -- (1161,111) -- cycle;
  \draw[PMist!60!white,line width=0.5pt,opacity=0.9] (1020,73) -- (1032,70) -- (1044,74) (1120,73) -- (1135,70) -- (1148,72);
  % buildings: thin, varied, low
  \foreach \x/\w/\h/\o in {958/3/6/.35,962/2/9/.4,966/3/7/.35,970/2/12/.45,974/3/9/.4,978/4/15/.45,983/2/11/.4,987/3/18/.5,
      991/2/13/.45,995/4/22/.55,1000/3/16/.5,1004/2/26/.55,1008/4/19/.5,1013/3/30/.6,1017/2/14/.45,1021/4/24/.55,1026/3/17/.5,
      1030/2/34/.6,1034/4/21/.55,1039/3/27/.55,1043/2/15/.45,1047/4/23/.55,1052/3/36/.6,1056/2/18/.5,1060/4/25/.55,1065/3/14/.45,
      1069/2/20/.5,1073/4/12/.45,1078/3/17/.5,1082/2/10/.4,1086/3/13/.45,1090/2/8/.4,1094/3/11/.4,1098/2/7/.35,1102/3/9/.35,1107/2/6/.35}{%
    \fill[PMist!85!PSub,opacity=\o] (\x,111) rectangle (\x+\w,111-\h);}
  \foreach \x/\h in {1031/34,1053/36,1005/26}{\draw[PMist!85!PSub,opacity=0.6,line width=0.5pt] (\x,111-\h) -- (\x,111-\h-7);}
  \fill[PRed!45!PMist,opacity=0.45] (1057,111) rectangle (1060,93);
  % pines at the far right
  \foreach \x/\h in {1124/12,1130/16,1136/13,1142/17,1148/14,1154/16,1159/12}{\Pine{\x}{111}{\h}{PTeal!70!black}}}


% ---------------------------------------------------------------- footer: left (grid + mountains) and right (Vancouver)
\newcommand{\FooterArtLeft}{%
  \begin{scope}
  \clip (12,1259) rectangle (77,1343);
  \fill[PTeal!40!PGreenDark,opacity=0.95] (22,1343) -- (52,1296) -- (68,1308) -- (96,1270) -- (118,1294) -- (134,1282) -- (176,1343) -- cycle;
  \draw[PTeal!35!white,line width=0.6pt,line join=round] (22,1343) -- (52,1296) -- (68,1308) -- (96,1270) -- (118,1294) -- (134,1282) -- (176,1343);
  \draw[white!80!PGreenDark,line width=0.6pt] (90,1278) -- (96,1283) -- (102,1278) (128,1290) -- (134,1287) -- (140,1291) (47,1303) -- (52,1300) -- (57,1304);
  \fill[PTeal!22!PGreenDark] (36,1343) -- (66,1320) -- (88,1331) -- (110,1318) -- (152,1343) -- cycle;
  \Pylon{22}{1341}{46}{white!82!PGreenDark}
  \Pylon{48}{1341}{32}{white!72!PGreenDark}
  \Pylon{66}{1341}{22}{white!62!PGreenDark}
  \draw[white!55!PGreenDark,line width=0.5pt] (14.6,1307) .. controls (30,1316) .. (42.9,1317.3) (29.4,1307) .. controls (41,1314) .. (53.1,1317.3)
     (51.1,1317.3) .. controls (58,1323) .. (63.6,1324.7);
  \foreach \x/\h in {172/18,180/26,188/20,196/24}{\Pine{\x}{1342}{\h}{PTeal!55!PGreenDark}}
  \end{scope}
  \draw[white!65!PGreenDark,line width=0.6pt] (79,1296) -- (79,1331);
  \node[anchor=west,text=white!90!PGreenDark,align=left,font=\CondSemi\fontsize{19pt}{24pt}\selectfont\addfontfeature{LetterSpace=6}] at (86,1313.5) {IEEE IAS ANNUAL MEETING\\OCT 4\textendash 8, 2026};}
\newcommand{\FooterArtRight}{%
  \begin{scope}
  \clip (970,1259) rectangle (1149,1343);
  \fill[PGreenDark!70!white,opacity=0.18] (985,1292) -- (1010,1276) -- (1024,1284) -- (1046,1265) -- (1070,1282) -- (1084,1275) -- (1106,1292) -- cycle;
  \draw[white!90!PGreenDark,line width=0.9pt,line join=round] (985,1292) -- (1010,1276) -- (1024,1284) -- (1046,1265) -- (1070,1282) -- (1084,1275) -- (1106,1292);
  \draw[white!90!PGreenDark,line width=0.6pt] (1040,1270) -- (1046,1274) -- (1052,1270) (1004,1280) -- (1010,1283) -- (1016,1280);
  \foreach \x/\h in {1116/24,1124/34,1133/28,1142/38,1151/30}{\Pine{\x}{1338}{\h}{white!90!PGreenDark}}
  \foreach \x/\h in {1109/16,1157/18}{\Pine{\x}{1339}{\h}{white!68!PGreenDark}}
  \draw[white!75!PGreenDark,line width=0.6pt] (1104,1339) -- (1159,1339);
  \end{scope}
  \node[anchor=west,text=white,font=\CondBold\fontsize{19.5pt}{19.5pt}\selectfont] at (984,1309) {VANCOUVER 2026};
  }
